\documentclass[12pt,a4paper]{article}
\usepackage[T1]{fontenc}
\usepackage{amsmath,amssymb,mathrsfs,tikz,times,pifont}
\usepackage{enumitem}
\usepackage{multicol}
\usepackage{lmodern}
\usetikzlibrary{trees}
\newcommand\circitem[1]{%
\tikz[baseline=(char.base)]{
\node[circle,draw=gray, fill=red!55,
minimum size=1.2em,inner sep=0] (char) {#1};}}
\newcommand\boxitem[1]{%
\tikz[baseline=(char.base)]{
\node[fill=cyan,
minimum size=1.2em,inner sep=0] (char) {#1};}}
\setlist[enumerate,1]{label=\protect\circitem{\arabic*}}
\setlist[enumerate,2]{label=\protect\boxitem{\alph*}}
%%%::::::by chnini ameur :::::::%%%
\everymath{\displaystyle}
\usepackage[left=1cm,right=1cm,top=1cm,bottom=1.7cm]{geometry}
\usepackage[colorlinks=true, linkcolor=blue, urlcolor=blue, citecolor=blue]{hyperref}
\usepackage{array,multirow}
\usepackage[most]{tcolorbox}
\usepackage{varwidth}
\usepackage{float} %pour utiliser l'option [H] qui force l'image à apparaître exactement à l'endroit où elle est placée dans le code.
\tcbuselibrary{skins,hooks}
\usetikzlibrary{patterns}
%%%::::::by chnini ameur :::::::%%%
\newtcolorbox{exa}[2][]{enhanced,breakable,before skip=2mm,after skip=5mm,
colback=yellow!20!white,colframe=black!20!blue,boxrule=0.5mm,
attach boxed title to top left ={xshift=0.6cm,yshift*=1mm-\tcboxedtitleheight},
fonttitle=\bfseries,
title={#2},#1,
% varwidth boxed title*=-3cm,
boxed title style={frame code={
\path[fill=tcbcolback!30!black]
([yshift=-1mm,xshift=-1mm]frame.north west)
arc[start angle=0,end angle=180,radius=1mm]
([yshift=-1mm,xshift=1mm]frame.north east)
arc[start angle=180,end angle=0,radius=1mm];
\path[left color=tcbcolback!60!black,right color = tcbcolback!60!black,
middle color = tcbcolback!80!black]
([xshift=-2mm]frame.north west) -- ([xshift=2mm]frame.north east)
[rounded corners=1mm]-- ([xshift=1mm,yshift=-1mm]frame.north east)
-- (frame.south east) -- (frame.south west)
-- ([xshift=-1mm,yshift=-1mm]frame.north west)
[sharp corners]-- cycle;
},interior engine=empty,
},interior style={top color=yellow!5}}
%%%%%%%%%%%%%%%%%%%%%%%
\usepackage{fancyhdr}
\usepackage{eso-pic}         % Pour ajouter des éléments en arrière-plan
% Commande pour ajouter du texte en arrière-plan
\usepackage{tkz-tab}
\AddToShipoutPicture{
    \AtTextCenter{%
        \makebox[0pt]{\rotatebox{80}{\textcolor[gray]{0.7}{\fontsize{5cm}{5cm}\selectfont PGB}}}
    }
}
\usepackage{lastpage}
\fancyhf{}
\pagestyle{fancy}
\renewcommand{\footrulewidth}{1pt}
\renewcommand{\headrulewidth}{0pt}
\renewcommand{\footruleskip}{10pt}
\fancyfoot[R]{
\color{blue}\ding{45}\ \textbf{2026}
}
\fancyfoot[L]{
\color{blue}\ding{45}\ \textbf{Prof:M. BA}
}
\cfoot{\bf
\thepage /
\pageref{LastPage}}
% Création du compteur pour les exercices
\newcounter{exercice}
\renewcommand{\theexercice}{\arabic{exercice}}  % Définit l'affichage du compteur en chiffres arabes

% Définir la commande \exo
\newcommand{\exo}{\refstepcounter{exercice}\textbf{Exercice \theexercice} }
\begin{document}
\renewcommand{\arraystretch}{1.5}
\renewcommand{\arrayrulewidth}{1.2pt}
\begin{tikzpicture}[overlay,remember picture]
    \node[draw=blue,line width=1.2pt,fill=purple,text=blue,inner sep=3mm,rounded corners,pattern=dots]at ([yshift=-2.5cm]current page.north) {\begingroup\setlength{\fboxsep}{0pt}\colorbox{white}{\begin{tabular}{|*1{>{\centering \arraybackslash}p{0.28\textwidth}} |*2{>{\centering \arraybackslash}p{0.2\textwidth}|} *1{>{\centering \arraybackslash}p{0.19\textwidth}|} }
                \hline
                \multicolumn{3}{|c|}{$\diamond$$\diamond$$\diamond$\ \textbf{Lycée de Dindéfélo}\ $\diamond$$\diamond$$\diamond$ } & \textbf{A.S. : 2026/2027}                                              \\ \hline
                \textbf{Matière: Mathématiques}                                                                                    & \textbf{Niveau : 2}\textbf{ndL} & \textbf{Date: 17/11/2026} & \textbf{} \\ \hline
                \multicolumn{4}{|c|}{\parbox[c]{10cm}{\begin{center}
                                                                  \textbf{{\Large\sffamily Td Calcul Dans $\mathbb{R}$}}
                                                              \end{center}}}                                                                                                        \\ \hline
            \end{tabular}}\endgroup};
\end{tikzpicture}

\vspace{3cm}
\fbox{\textbf{\exo}}\textbf{Date: 19/11/2026}

Ecrire les réels suivants sous formes de fractions irréductibles :

\begin{multicols}{2}
\setlength{\columnseprule}{0.1mm}
\begin{enumerate}
\item .

$\begin{aligned}
            \textbf{A} &= \frac{\frac{3}{2}-\frac{5}{3}}{-\frac{7}{2}+\frac{1}{6}} \times \frac{1}{-3+\frac{2}{5}} \\[0.8em]
            &= \frac{\frac{9}{6}-\frac{10}{6}}{-\frac{21}{6}+\frac{1}{6}} \times \frac{1}{-\frac{15}{5}+\frac{2}{5}} \\[0.8em]
            &= \frac{-\frac{1}{6}}{-\frac{20}{6}} \times \frac{1}{-\frac{13}{5}} \\[0.8em]
            &= \left(-\frac{1}{6}\right) \times \left(-\frac{6}{20}\right) \times \left(-\frac{5}{13}\right) \\[0.5em]
            &= -\frac{1 \times 6 \times 5}{6 \times 20 \times 13} \\[0.5em]
            &= -\frac{5}{20 \times 13} \\[0.5em]
            &= -\frac{1}{4 \times 13} \\[0.5em]
            &= -\frac{1}{52}
        \end{aligned}
$
        
\item .

$\begin{aligned}
            \textbf{B} &= \frac{2}{3}-\frac{4}{5}+\frac{3}{7} \\[0.5em]
            &= \frac{2 \times 35}{3 \times 35} - \frac{4 \times 21}{5 \times 21} + \frac{3 \times 15}{7 \times 15} \\[0.5em]
            &= \frac{70}{105} - \frac{84}{105} + \frac{45}{105} \\[0.5em]
            &= \frac{70 - 84 + 45}{105} \\[0.5em]
            &= \frac{-14 + 45}{105} \\[0.5em]
            &= \frac{31}{105}
        \end{aligned}$
        
\begin{comment}
\item $\textbf{C}=\left(3+\frac{4}{5} \right) \left(\frac{1}{5}-4 \right)$
\item $\textbf{D}=\left(\frac{2}{3}-1 \right)\times \frac{2+\frac{1}{3}}{\frac{3}{5}+\frac{1}{5}}$
\item $\textbf{E}=\frac{3+\frac{1}{3}}{\frac{2}{3}-2} \times \frac{2}{\frac{1}{3}+3}$
\item $\textbf{F}=\frac{\frac{1}{4}+3}{\frac{3}{5}} \times \frac{7}{3}$
\item $\textbf{G}= \frac{5}{2}-\frac{1}{\frac{2}{3}}$
\item $\textbf{H}= \frac{\frac{2}{3}-\frac{2}{3}}{\frac{5}{4}} - \frac{7}{15}$
\item $\textbf{I}= \frac{\frac{1}{2}+\frac{2}{3}-1}{\frac{3}{4}}$
\item $\textbf{J}= \left(5-\frac{3}{2} \right)  \left(\frac{3}{4}+\frac{-5}{4} \right)$
\item $\textbf{K}=\left(\frac{5}{3}+1 \right) \div \left(\frac{4}{3}-\frac{1}{4} \right)  $
\item $\textbf{L}=\left(3-\frac{2}{3} \right) + \left(-\frac{3}{2}+3 \right)   $
\end{comment}
\end{enumerate}
\end{multicols}

\end{document}